\section{Related Work}

\paragraph{Memory for LLM Agents.}
Memory has become a foundational capability for LLM agents, supporting long-horizon understanding, continual adaptation in complex environments\cite{zhang2024survey_memoery,wu2025humansurveymemory,hu2025memoryageaiagentssurvey}.
To equip LLM agents with memory, most methods construct an explicit memory bank supported by primitives for  segmentation \cite{secom}, summarization \cite{kim2024theanine,RecurSum,SumMem,Think-in-memory,memochat,rasmussen2025zep}, compression \cite{chen2025compress,ReadAgent,recomp,COMEDY}, and forgetting/updating to maintain long-term quality \cite{memorybank,LD-Agent}.
To improve retrieval quality, several approaches build structured memory indices such as trees \cite{memtree,raptor} and graphs \cite{hipporag2,chhikara2025mem0,wang2025mirix,xu-etal-2024-multi}. Beyond generic storage, personalization-oriented methods emphasize capturing user profiles and preferences for downstream conditioning \cite{zhang2026personalize,xu2025harnessing,zhang2025personaagent,perltqa,fang2025lightmem,memorag}. Another thread focuses on experience memory, where agents memorize successful or failed trajectories to improve future decision-making \cite{packer2023memgpt,AWM,tang2025agentkb,ouyang2025reasoningbank,zhao2024expel,zhang2025gmemory,gao2025llm4rerank,jia2024mill}.  Despite strong performance, these methods rely on hand-crafted heuristics, which can be brittle under non-stationary user behaviors. Our method is orthogonal to memory construction and storage, and is broadly compatible with existing memory backends as a general mechanism for memory evolution to support long-term personalization.
% In contrast, our approach is orthogonal to specific memory construction and storage protocols; it serves as a pluggable mechanism for memory evolution, offering broad compatibility with existing memory backends to support long-term personalization.

\paragraph{RL for Memory.}
Recent works treat memory operations as a sequential decision problem and optimize it with reinforcement learning~\cite{zhang2026evoking,wang2025memalpha,zhou2025mem1,memory-r1,m3-agent,yuan2025memsearcher,zhang2025MemoryasAction,liu2025large,li2025towards}.
For example, RMM~\cite{RMM} learns to manage long-term personalized memory via reflective update and retrieval; MemAgent~\cite{yu2025memagent} uses RL to learn a memory agent that maintains a fixed-length context by selectively preserving/overwriting long dialogue history; MEM1~\cite{zhou2025mem1} trains memory and reasoning synergy to form compact memory for efficient long-horizon agent; MemGen~\cite{zhang2025memgen} proposes generative latent memory that weaves experience into reusable memory tokens for self-evolving agents.
However, these approaches typically rely on final success/answer as sparse rewards, and lack process-level rewards that directly guide how memory should be updated. We address this by introducing \emph{guideline-aligned rewards}, which provide structured learning signals for memory evolution.

\paragraph{Prompt Optimization.}
A growing line of work treats prompts as optimizable natural-language parameters, iteratively refining them via model-generated feedback rather than numerical gradients~\cite{yuksekgonul2025optimizing,yang2023large,pryzant2023automatic,shinn2023reflexion,tang2025unleashing,zhang2024offline,zhang2024revolve,zhang2024aflow}.
The common pattern involves evaluating the current prompt, generating natural-language edit signals (textual gradients), and applying them to produce improved variants.
For instance, TextGrad~\cite{yuksekgonul2025optimizing} uses LLM-generated textual gradients for iterative prompt refinement; OPRO~\cite{yang2023large} treats the LLM as a black-box optimizer that proposes and scores prompt candidates; Reflexion~\cite{shinn2023reflexion} converts past failures into self-reflection feedback carried forward to guide future actions.
While related in spirit, MGI differs in that it optimizes a global memory-evolution guideline over long histories rather than a single-turn prompt, stabilizes updates via contrastive diagnosis and batch aggregation, and interfaces with Stage-2 RL through guideline-aligned process rewards to jointly optimize memory update policies and reinforced behaviors.